\scanpage{26}
\mspage{\unsure{35|55}}
tel que $A$ soit séparé et complet pour la topologie $J$-adique. Sous ces conditions:

(i) Le foncteur $F\rightsquigarrow\widehat F$ de la catégorie des Modules localement libres
sur $U'$ dans la catégorie des Modules localement libres sur $\widehat U'$ est
pleinement fidèle.

(ii) Pour tout Module localement libre $\underline F$ sur $\widehat U'$, il existe un Module
cohérent $F$ sur $\unsure{U'|V'}$ (pas nécessairement localement libre, hélas), et un
isomorphisme $\widehat F\simeq\underline F$.

Il reste seulement à prouver (ii), grâce à 5.9. Or d'après \unsure{ceytte|cette}
remarque, et exposé IX 2.1. il s'ensuit que $\underline F$ est induit par un Module
cohérent $\underline G$ sur $\widehat X'$. D'après le théorème \deleted{de comparaison et} d'existence\inserted{EGA $\overline{\underline{\mathrm{III}}}$}
5.1.4., $\underline G$ est de la forme $\widehat F$, où $F$ est cohérent sur $X$, d'où la conclusion.

\uline{Remarques} 5.1\unsure{0|1}.\inserted{$1^{\mathrm o})$} Utilisant 5.10.\inserted{, 4.7.}
\deleted{et des hypothèses convenables de}
une hypothèse convenable, disant que certains anneaux locaux des fibres
géométriques de $X'\quad S'$ sont ``purs'' resp. \unsure{parafactoriels|Parafactoriels}, on doit pouvoir
obtenir des énoncés disant \deleted{\unsure{que}} que le foncteur $Z'\rightsquigarrow\widehat Z'$ de la catégorie
des revêtements étales de $X'$ dans la catégorie des revêtements étales de
$\widehat X'$ \deleted{est} (ou ce qui revient au même, de $X'_0$) est \deleted{pleinement fidèle} une
équivalence de catégories, resp. que le foncteur $\underline L\rightsquigarrow\widehat{\underline L}$ de la catégorie
des Modules inversibles sur $X'$ dans la catégorie des Modules inversibles
sur $\widehat X'$ est une équivalence. Utilisant des résultats récents de MURRE,
il est probable que l'on doit pouvoir en déduire des théorèmes d'existence
\deleted{pour} des schémas de Picard pour \deleted{des}\inserted{certains} schémas algébriques \deleted{propres} \uline{non}
\uline{propres}. \deleted{\unsure{voir}} De façon générale, l'élimination d'hypothèses de propreté
dans divers théorèmes d'existence, notamment de représentabilité de
foncteurs comme les foncteurs de Hilbert, ou de Picard etc, à l'aide
\margin{25}

